%% kapa: exam=1 round=20 year=2009 period=2 form=A subject=appraisal_law subject_ko=부동산관계법규 seq=26 number=66
%% source: Q-Net 감정평가사 기출문제 (cst00309) artlSeq=5024178 file=extracted/1008/5024178/2009년 제20회 감정평가사 1차시험 기출문제/2교시.hwp
%% conversion: kordoc 4.24.51 -> tools/gen1.py (auto); review=proofread-sonnet; flags=none
\begin{question}{66}
국유재산법령상의 내용으로 옳지 \uline{않은} 것은?
\begin{choices}
  \item 국유재산에 관한 사무를 총괄하는 총괄청은 기획재정부장관이다.
  \item 총괄청 및 관리청은 국유재산의 범위에 속하는 재산의 기부에 조건이 수반된 것인 경우에도 이를 채납할 수 있다.
  \item 행정재산에 대해서는 사권을 설정할 수 없다.
  \item 행정재산이라도 직접 공용으로 사용하기 위하여 필요로 하는 지방자치단체에 양여하는 것은 가능하다.
  \item 국가가 장차 6년 후에 직접 공공용으로 사용하기로 결정한 재산은 행정재산에 해당하지 않는다.
\end{choices}
\end{question}
\answer{2}
